\documentclass[10pt,a4paper]{amsart}
\usepackage[margin=19mm]{geometry}
\usepackage{amsmath,amssymb,mathtools}
\usepackage[scr=rsfs,cal=euler]{mathalfa}
\usepackage{xfrac}
\usepackage[unicode,hidelinks,bookmarks=false]{hyperref}
\newcommand{\ie}{i.e.\ }
\newcommand{\cf}{cf.\ }
\newcommand{\resp}{respectively\ }
\newcommand{\Subsection}{\subsection}
\providecommand{\nobreakdash}{\nobreak\hbox{-}}
\setlength{\emergencystretch}{2em}
\multlinegap0pt
\usepackage{amssymb}
\newtheorem{lemma}{Lemma}
\DeclareMathOperator{\Aut}{Aut}
\newcommand*{\C}{{\mathbb C}}
\newcommand*{\R}{{\mathbb R}}
\DeclareMathOperator{\pr}{pr}
\title{On holomorphic submersions with biholomorphic fibers}
\author{Serge Lvovski}
\date{Source: arXiv:2608.08004v2 (CC BY 4.0)}
\begin{document}
\maketitle

\noindent\emph{Source and licence.} On holomorphic submersions with biholomorphic fibers. Version arXiv:2608.08004v2 (CC BY 4.0), distributed by arXiv under the Creative Commons Attribution 4.0 International licence, http://creativecommons.org/licenses/by/4.0/, as recorded in the arXiv licence metadata for this exact version and shown on the abstract page https://arxiv.org/abs/2608.08004v2. Full licence notice: \texttt{licenses/arxiv-2608.08004-CC-BY-4.0.txt}.\par\smallskip

\noindent\textit{Editorial scope.} Counted unit: the article's lemma lemma (source0.tex lines 547--558). The article's complete original proof of it is delivered with it (source0.tex lines 560--611, one \texttt{proof} environment of 52 lines). No mathematics is written, repaired or re-derived: the statement and the proof are the article's own lines, in the article's own notation, and no retained line is substituted or re-flowed.  No further source-local context is needed: every cross-reference of every retained line resolves inside the two delivered spans.  Nothing was omitted from any retained span, so the package carries no omission record.  No retained line contains a citation, so the wrapper carries no bibliography and no bibliography entry.  No retained line contains a figure environment, an image-inclusion command or drawing code, so no image is delivered and the package declares no asset dependency.  Editorial formatting only, in addition: an AMS \texttt{amsart} preamble that restates the article's own declarations and macros used by the retained lines, namely the environments \texttt{lemma} on separate counters, together with the standard packages those lines need.  The retained environments print in this document as Lemma 1 in order of appearance, whereas the article prints the same items with the numbering of its own full document.  The complete native source of the article as distributed in the arXiv e-print for this exact version is bundled unchanged as \texttt{sources/207123/source0.tex}.
\begin{lemma}\label{lemma}
Suppose that $V\subset\C^n$ is an open subset and the mapping $g\colon
V\to\Aut(D)$, where $D\subset\C$ is the unit disc, is such that 
\begin{equation}\label{eq:G}
\begin{aligned}
  G\colon &V\times D\to V\times D,\\
  &(z,w)\mapsto(z,g(z)w)  
\end{aligned}
\end{equation}
is a biholomorphism from $V\times D$ to
itself. Then $g$ is locally constant.
\end{lemma}

\begin{proof}
For any $z\in V$ put $s(z)=\pr_2(G^{-1}(z,0))$; since $G$ is holomorphic,
the mapping $z\mapsto s(z)$ is a holomorphic mapping from $V$
to~$D$. Hence, $V$ can be covered by open subspaces~$V_\alpha$ such
that, for each~$\alpha$,
\begin{equation}\label{eq:lemma}
  G|_{V_\alpha\times D}\colon (z,w)\mapsto e^{i\theta(z)}
  \frac{w-s(z)}{1-\overline{s(z)}w},  
\end{equation}
where $\theta(z)\in\R$; without loss of generality we may and will
assume that $V=V_\alpha$
for some~$\alpha$ and that
$V$ is open and connected. If $z\in V\subset\C^n$, put $z=(z_1,\dots,z_n)$.

Since the right-hand side of~\eqref{eq:lemma} is holomorphic as a
function of~$z$, on has, for any fixed~$w\in D$ and for any~$j$, $1\le
j\le n$,
\begin{multline*}
  0\equiv\frac\partial{\partial \bar z_j}
  \left(e^{i\theta(z)}\frac{w-s(z)}{1-\overline{s(z)}w}\right)\\
{}=e^{i\theta(z)}\left(
\frac{w-s(z)}{1-\overline{s(z)}w}\cdot
i\frac{\partial\theta(z)}{\partial \bar z_j}+
\frac{w-s(z)}{(1-\overline{ s(z)}w)^2}\cdot
\frac{\partial\overline{s(z)}}{\partial \bar z_j}\right)
\end{multline*}
for any $w\in D$ and any $j$, whence
\begin{equation*}
  i(1-\overline{s(z)}w)\frac{\partial\theta(z)}{\partial \bar z_j}
  + \overline{\left(\frac{\partial s}{\partial z_j}\right)}
  =0\quad\text{for
any $w\in D$ and any $j$,}  
\end{equation*}
or, equivalently,
\begin{align}
  i\frac{\partial\theta(z)}{\partial \bar
    z_j}+\overline{\left(\frac{\partial s}{\partial z_j}\right)}
  &=0,\label{eqa}\\ 
\overline{s(z)}\cdot \frac{\partial\theta(z)}{\partial \bar
    z_j}&=0\label{eqb}  
\end{align}
for any~$j$, $1\le j\le n$.

If the function~$s$ is not identically zero, then it follows
from~\eqref{eqb} that $\theta$ is holomorphic; being real valued, it
must be constant. Now~\eqref{eqa} implies that $s$ is constant, too,
and we are done.

If the function~$s$ is identically zero, then \eqref{eqa} implies that
the real valued function~$\theta$ is holomorphic, hence 
constant, and we are done.
\end{proof}

\end{document}
